\begin{textsample}{POS Dim 3 – ap – Score 14.00 – ap/ap\_inf\_11.txt}  \label{ex:f3_pos_146}
III – Campo de Experiência : Traços, \textbf{Sons} Cores e Formas : Conviver com diferentes manifestações artísticas, culturais e científicas, locais e universais, no cotidiano da instituição escolar, possibilita às \textbf{crianças}, por meio de experiências diversificadas, vivenciar diversas formas de expressão e linguagens, como as artes visuais ( \textbf{pintura}, modelagem, colagem, fotografia etc. ), a música, o teatro, a dança e o audiovisual, entre outras.
Com base nessas experiências, elas se expressam por várias linguagens, criando suas próprias produções artísticas ou culturais, exercitando a autoria ( coletiva e individual ) com \textbf{sons}, traços, \textbf{gestos}, danças, \textbf{mímicas}, encenações, canções, desenhos, modelagens, manipulação de diversos materiais e de recursos tecnológicos.
Essas experiências contribuem para que, desde muito \textbf{pequenas}, as \textbf{crianças} desenvolvam senso estético e crítico, o conhecimento de si mesmas, dos outros e da realidade que as cerca.
Portanto, a Educação \textbf{Infantil} precisa promover a participação das \textbf{crianças} em tempos e espaços para a produção, manifestação e apreciação artística, de modo a favorecer o desenvolvimento da sensibilidade, da criatividade e da expressão pessoal das \textbf{crianças}, permitindo que se apropriem e reconfigurem, permanentemente, a cultura e potencializem suas singularidades, ao ampliar repertórios e interpretar suas experiências e vivências artísticas.
Direitos de aprendizagem com foco no campo de experiência traços, \textbf{sons}, cores e formas : - CONVIVER e fruir das manifestações artísticas e culturais da sua comunidade e de outras culturas – artes \textbf{plásticas}, música, dança, teatro, cinema, folguedos e festas populares - ampliando a sua sensibilidade, desenvolvendo senso estético, empatia e respeito às diferentes culturas e identidades. - \textbf{BRINCAR} com diferentes \textbf{sons}, ritmos, formas, cores, \textbf{texturas}, objetos, materiais, construindo cenários e indumentárias para \textbf{brincadeiras} de faz de conta, encenações ou para festas tradicionais, enriquecendo seu repertório e desenvolvendo seu senso estético. - PARTICIPAR de decisões e ações relativas à organização do ambiente ( tanto no cotidiano como na preparação de eventos especiais ), à definição de temas e à escolha de materiais a serem usados em atividades \textbf{lúdicas} e teatrais, entrando em contato com manifestações do patrimônio cultural, artístico e tecnológico, apropriando -se de diferentes linguagens. - EXPLORAR variadas possibilidades de usos e combinações de materiais, substâncias, objetos e recursos tecnológicos para criar e recriar danças, artes visuais, encenações teatrais, músicas, escritas e mapas, apropriando -se de diferentes manifestações artísticas e culturais. - EXPRESSAR, com criatividade e responsabilidade, suas emoções, sentimentos, necessidades e ideias \textbf{brincando}, cantando, dançando, esculpindo, desenhando, encenando, compreendendo e usufruindo ( do ) o que é comunicado pelos demais \textbf{colegas} e pelos \textbf{adultos}. - CONHECER -SE, no contato criativo com manifestações artísticas e culturais locais e de outras comunidades, identificando e valorizando o seu pertencimento étnico-racial, de gênero e de crença religiosa, desenvolvendo sua sensibilidade, criatividade, gosto pessoal e modo peculiar de expressão por meio do teatro, música, dança, desenho e imagens.

% matched lemmas: adulto, brincadeira, brincar, colega, criança, gesto, infantil, lúdico, mímico, pequeno, pintura, plástico, som, textura
\end{textsample}
